\maketitle

\begin{abstract}
I present an open, reproducible firm-level microsimulation for costing UK VAT registration-threshold reforms, released with its synthetic-data generator. The threshold is a \emph{notch}: once annual turnover crosses \pounds85{,}000, VAT falls due on a firm's whole base, not the excess. I build synthetic firm records from ONS registered-business counts, calibrate them to 2023--24 HMRC statistics, and cost \emph{level}, \emph{shape} (a taper) and \emph{rate} (a reduced-rate band) reforms statically on a common \pounds184.7bn base: raising the threshold to \pounds100{,}000 costs \pounds698m, a graduated taper over \pounds85{,}000--\pounds105{,}000 costs \pounds466m, and 10\% and 15\% reduced-rate bands cost \pounds460m and \pounds230m. A static costing of the April 2024 rise to \pounds90{,}000 brackets HMRC's published costing between full-deregistration and voluntary-retention conventions in every forecast year. I characterise the notch by its \emph{dominated region}---the exact \pounds21{,}250 band above the threshold in which no firm profits from locating---and how each reform changes it: a level rise relocates it, a reduced rate splits it into two bands of essentially unchanged total width, a taper removes it. Conditional on an assumed turnover elasticity, intensive-margin responses barely move these costings: a pure level rise has exactly zero intensive-margin offset, and reduced-rate offsets stay under 4\%. The synthetic data contain no behavioural bunching---the estimator finds none at \pounds85{,}000 and a large spurious signal at a \pounds90{,}000 calibration-band edge---so the paper's transparency result is that band-calibrated synthetic data cannot support bunching inference in either direction.
\end{abstract}

\vspace{0.3cm}
\noindent\textbf{Keywords:} value-added tax, microsimulation, tax notch, registration threshold

\newpage
